\documentclass[10pt,a4paper]{amsart}
\usepackage{amsmath,amssymb,amsthm,graphicx}
% Retained preamble for the Stacks Project chapter {CH} (shared project preamble).
% The chapter file pinned-original-source.tex only carries \input{{preamble}} on its first line,
% so no preamble text lives inside the chapter itself.  The Stacks Project uses one single shared
% preamble for every chapter: the snapshot stages it three times (sources/stacks-chow/preamble.tex,
% sources/stacks-etale-cohomology/preamble.tex and focused-seeds/AGP/sources/stacks-snapshot-20261011/
% preamble.tex) and those copies are identical in size.  The block below is copied mechanically,
% line for line, from sources/stacks-chow/preamble.tex lines 14-18 and 154-217: the Xy-pic setup and
% the native theorem declarations together with the mathematical macros.  Nothing was added, removed
% or redefined; the wrapper class and the amsmath/amssymb/amsthm/graphicx base packages are supplied
% by the assembler.
% For commutative diagrams we use Xy-pic
% Chapter: stacks-cohomology.tex -- Cohomology of Sheaves.
\usepackage[all]{xy}

% We use 2cell for 2-commutative diagrams.
\xyoption{2cell}

% Theorem environments.
%
\theoremstyle{plain}
\newtheorem{theorem}[subsection]{Theorem}
\newtheorem{proposition}[subsection]{Proposition}
\newtheorem{lemma}[subsection]{Lemma}

\theoremstyle{definition}
\newtheorem{definition}[subsection]{Definition}
\newtheorem{example}[subsection]{Example}
\newtheorem{exercise}[subsection]{Exercise}
\newtheorem{situation}[subsection]{Situation}

\theoremstyle{remark}
\newtheorem{remark}[subsection]{Remark}
\newtheorem{remarks}[subsection]{Remarks}

\numberwithin{equation}{subsection}

% Macros
%
\def\lim{\mathop{\mathrm{lim}}\nolimits}
\def\colim{\mathop{\mathrm{colim}}\nolimits}
\def\Spec{\mathop{\mathrm{Spec}}}
\def\Hom{\mathop{\mathrm{Hom}}\nolimits}
\def\Ext{\mathop{\mathrm{Ext}}\nolimits}
\def\SheafHom{\mathop{\mathcal{H}\!\mathit{om}}\nolimits}
\def\SheafExt{\mathop{\mathcal{E}\!\mathit{xt}}\nolimits}
\def\Sch{\mathit{Sch}}
\def\Mor{\mathop{\mathrm{Mor}}\nolimits}
\def\Ob{\mathop{\mathrm{Ob}}\nolimits}
\def\Sh{\mathop{\mathit{Sh}}\nolimits}
\def\NL{\mathop{N\!L}\nolimits}
\def\CH{\mathop{\mathrm{CH}}\nolimits}
\def\proetale{{pro\text{-}\acute{e}tale}}
\def\etale{{\acute{e}tale}}
\def\QCoh{\mathit{QCoh}}
\def\Ker{\mathop{\mathrm{Ker}}}
\def\Im{\mathop{\mathrm{Im}}}
\def\Coker{\mathop{\mathrm{Coker}}}
\def\Coim{\mathop{\mathrm{Coim}}}

% Boxtimes
%
\DeclareMathSymbol{\boxtimes}{\mathbin}{AMSa}{"02}

%
% Macros for moduli stacks/spaces
%
\def\QCohstack{\mathcal{QC}\!\mathit{oh}}
\def\Cohstack{\mathcal{C}\!\mathit{oh}}
\def\Spacesstack{\mathcal{S}\!\mathit{paces}}
\def\Quotfunctor{\mathrm{Quot}}
\def\Hilbfunctor{\mathrm{Hilb}}
\def\Curvesstack{\mathcal{C}\!\mathit{urves}}
\def\Polarizedstack{\mathcal{P}\!\mathit{olarized}}
\def\Complexesstack{\mathcal{C}\!\mathit{omplexes}}
% \Pic is the operator that assigns to X its picard group, usage \Pic(X)
% \Picardstack_{X/B} denotes the Picard stack of X over B
% \Picardfunctor_{X/B} denotes the Picard functor of X over B
\def\Pic{\mathop{\mathrm{Pic}}\nolimits}
\def\Picardstack{\mathcal{P}\!\mathit{ic}}
\def\Picardfunctor{\mathrm{Pic}}
\def\Deformationcategory{\mathcal{D}\!\mathit{ef}}

% --- appended 2026-10-12 (Chief): the project own definitions of the reference/slogan/history
% display environments, copied mechanically from the shared Stacks preamble lines 8-12
% (\usepackage{verbatim} plus the three \newenvironment lines, which the chapter preamble keeps
% verbatim from the project preamble). Nothing added, removed or redefined: these environments
% carry no mathematical content and the project itself typesets them as comments. Without these
% three lines a counted statement whose own text uses them fails with Environment slogan undefined.
\usepackage{verbatim}
\newenvironment{reference}{\comment}{\endcomment}
\newenvironment{slogan}{\comment}{\endcomment}
\newenvironment{history}{\comment}{\endcomment}
\usepackage{cleveref}
\crefname{originallemma}{original lemma}{original lemmas}
\Crefname{originallemma}{Original lemma}{Original lemmas}
\makeatletter
\newlabel{lemma-quasi-separated-cohomology-colimit}{{01FF}{source}{}{sourcepointer-1}{}}
\newlabel{lemma-quasi-separated-cohomology-colimit@cref}{{[originallemma][0][]01FF}{[1][0][]source}{}{}{}}
\makeatother
\crefname{originallemma}{original lemma}{original lemmas}
\Crefname{originallemma}{Original lemma}{Original lemmas}
\makeatletter
\newlabel{lemma-subsheaf-of-constant-sheaf}{{0A38}{source}{}{sourcepointer-2}{}}
\newlabel{lemma-subsheaf-of-constant-sheaf@cref}{{[originallemma][0][]0A38}{[1][0][]source}{}{}{}}
\makeatother
\crefname{originaldefinition}{original definition}{original definitions}
\Crefname{originaldefinition}{Original definition}{Original definitions}
\makeatletter
\newlabel{modules-definition-generated-by-local-sections}{{01AQ}{source}{}{sourcepointer-3}{}}
\newlabel{modules-definition-generated-by-local-sections@cref}{{[originaldefinition][0][]01AQ}{[1][0][]source}{}{}{}}
\makeatother
\crefname{originallemma}{original lemma}{original lemmas}
\Crefname{originallemma}{Original lemma}{Original lemmas}
\makeatletter
\newlabel{modules-lemma-generated-by-local-sections-stalk}{{01AR}{source}{}{sourcepointer-4}{}}
\newlabel{modules-lemma-generated-by-local-sections-stalk@cref}{{[originallemma][0][]01AR}{[1][0][]source}{}{}{}}
\makeatother
\title{Vanishing above the dimension from the vanishing on extensions by zero of basis opens}
\author{The Stacks Project Authors}
\date{stacks-project:cohomology.tex:local-snapshot-20261011; GNU FDL 1.2 or later}
\begin{document}
\maketitle
\noindent\textit{Editorial note.} All delivered mathematical text is the author's own, copied from the pinned Stacks Project chapter file cohomology.tex; no statement or proof step was rewritten, reordered or removed.\par\medskip
\noindent\textit{Editorial note.} This card was re-admitted after a gate audit: it had been excluded only because the chapter's retained preamble previously lacked the reference/slogan/history environments. The shared preamble now carries them (appended 2026-10-12), so the card uses that file as its preamble with no per-card copy.\par\medskip
\noindent\textit{Editorial note.} The author's own bibliographic remark, which sits in a reference environment at the head of the statement and describes the result as a special case of a proposition of the cited paper, is preserved in the source span and is not typeset because the project's own preamble defines reference environments as comments; no bibliography excerpt is needed.\par\medskip
\noindent\textit{Editorial note.} The label lemma-vanishing-generated-one-section is defined only in cohomology.tex, so no cross-chapter restatement had to be weighed for this card.\par\medskip
\noindent\textit{Allowed standard prerequisites (editorial).}\begin{itemize}
\item Any sheaf is the filtered colimit of the subsheaves generated by finitely many local sections. Modules, Lemma modules-lemma-generated-by-local-sections-stalk and Definition modules-definition-generated-by-local-sections, invoked at proof lines 3550-3557. Supplies the first reduction to finitely generated sheaves.
\item Cohomology commutes with filtered colimits of sheaves on a quasi-separated space. Same chapter, Lemma lemma-quasi-separated-cohomology-colimit, invoked at proof lines 3562-3567. Supplies the passage from the colimit of sheaves to the colimit of cohomology.
\item Filtration of a finitely generated subsheaf of the constant sheaf Z by extensions of extensions by zero. Same chapter, Lemma lemma-subsheaf-of-constant-sheaf, invoked at proof lines 3601-3610. Supplies the filtration that transfers the hypothesis to the successive quotients.
\end{itemize}\medskip
\ifdefined\hypertarget\hypertarget{sourcepointer-1}{}\fi
\noindent\textit{Source reference.} Original lemma 01FF: sources/stacks-cohomology/pinned-original-source.tex, source lines 3194--3194. Pointer outside the counted span, invoked once for context; the load-bearing argument of the counted statement is the author's proof reproduced in full in the proof span.\par\medskip
\ifdefined\hypertarget\hypertarget{sourcepointer-2}{}\fi
\noindent\textit{Source reference.} Original lemma 0A38: sources/stacks-cohomology/pinned-original-source.tex, source lines 3480--3480. Pointer outside the counted span, invoked once for context; the load-bearing argument of the counted statement is the author's proof reproduced in full in the proof span.\par\medskip
\ifdefined\hypertarget\hypertarget{sourcepointer-3}{}\fi
\noindent\textit{Source reference.} Original definition 01AQ: sources/stacks-chow/tags.txt, source lines 1699--1699. Pointer outside the counted span, invoked once for context; the load-bearing argument of the counted statement is the author's proof reproduced in full in the proof span.\par\medskip
\ifdefined\hypertarget\hypertarget{sourcepointer-4}{}\fi
\noindent\textit{Source reference.} Original lemma 01AR: sources/stacks-chow/tags.txt, source lines 1700--1700. Pointer outside the counted span, invoked once for context; the load-bearing argument of the counted statement is the author's proof reproduced in full in the proof span.\par\medskip
\begin{lemma}
\label{lemma-vanishing-generated-one-section}
\begin{reference}
This is a special case of \cite[Proposition 3.6.1]{Tohoku}.
\end{reference}
Let $X$ be a topological space. Let $d \geq 0$ be an integer. Assume
\begin{enumerate}
\item $X$ is quasi-compact,
\item the quasi-compact opens form a basis for $X$, and
\item the intersection of two quasi-compact opens is quasi-compact.
\item $H^p(X, j_!\underline{\mathbf{Z}}_U) = 0$ for all $p > d$
and any quasi-compact open $j : U \to X$.
\end{enumerate}
Then $H^p(X, \mathcal{F}) = 0$ for all $p > d$
and any abelian sheaf $\mathcal{F}$ on $X$.
\end{lemma}

\begin{proof}
Let $S = \coprod_{U \subset X} \mathcal{F}(U)$ where $U$ runs over the
quasi-compact opens of $X$.
For any finite subset $A = \{s_1, \ldots, s_n\} \subset S$,
let $\mathcal{F}_A$ be the subsheaf of $\mathcal{F}$ generated
by all $s_i$ (see
Modules, Definition \ref{modules-definition-generated-by-local-sections}).
Note that if $A \subset A'$, then $\mathcal{F}_A \subset \mathcal{F}_{A'}$.
Hence $\{\mathcal{F}_A\}$ forms a system over the
directed partially ordered set of finite subsets of $S$.
By Modules, Lemma \ref{modules-lemma-generated-by-local-sections-stalk}
it is clear that
$$
\colim_A \mathcal{F}_A = \mathcal{F}
$$
by looking at stalks. By
Lemma \ref{lemma-quasi-separated-cohomology-colimit} we have
$$
H^p(X, \mathcal{F}) =
\colim_A H^p(X, \mathcal{F}_A)
$$
Hence it suffices to prove the vanishing for the abelian sheaves
$\mathcal{F}_A$. In other words, it suffices to prove the
result when $\mathcal{F}$ is generated by finitely many local sections
over quasi-compact opens of $X$.

\medskip\noindent
Suppose that $\mathcal{F}$ is generated by the local sections
$s_1, \ldots, s_n$. Let $\mathcal{F}' \subset \mathcal{F}$
be the subsheaf generated by $s_1, \ldots, s_{n - 1}$.
Then we have a short exact sequence
$$
0 \to \mathcal{F}' \to \mathcal{F} \to \mathcal{F}/\mathcal{F}' \to 0
$$
From the long exact sequence of cohomology we see that it suffices
to prove the vanishing for the abelian sheaves $\mathcal{F}'$
and $\mathcal{F}/\mathcal{F}'$ which are generated by fewer than
$n$ local sections. Hence it suffices to prove the vanishing
for sheaves generated by at most one local section. These sheaves
are exactly the quotients of the sheaves $j_!\underline{\mathbf{Z}}_U$
where $U$ is a quasi-compact open of $X$.

\medskip\noindent
Assume now that we have a short exact sequence
$$
0 \to \mathcal{K} \to j_!\underline{\mathbf{Z}}_U \to \mathcal{F} \to 0
$$
with $U$ quasi-compact open in $X$.
It suffices to show that $H^q(X, \mathcal{K})$ is zero for $q \geq d + 1$.
As above we can write $\mathcal{K}$ as the filtered colimit of
subsheaves $\mathcal{K}'$ generated by finitely many sections over
quasi-compact opens. Then $\mathcal{F}$ is the filtered colimit of the
sheaves $j_!\underline{\mathbf{Z}}_U/\mathcal{K}'$. In this way we
reduce to the case that $\mathcal{K}$ is generated by finitely many
sections over quasi-compact opens. Note that $\mathcal{K}$
is a subsheaf of $\underline{\mathbf{Z}}_X$. Thus by
Lemma \ref{lemma-subsheaf-of-constant-sheaf} there exists a finite
filtration of $\mathcal{K}$ whose successive quotients $\mathcal{Q}$ fit
into a short exact sequence
$$
0 \to j''_!\underline{\mathbf{Z}}_W \to
j'_!\underline{\mathbf{Z}}_V \to \mathcal{Q} \to 0
$$
with $j'' : W \to X$ and $j' : V \to X$ the inclusions of quasi-compact opens.
Hence the vanishing of $H^p(X, \mathcal{Q})$ for $p > d$ follows
from our assumption (in the lemma) on the vanishing of the cohomology groups
of $j''_!\underline{\mathbf{Z}}_W$ and $j'_!\underline{\mathbf{Z}}_V$.
Returning to $\mathcal{K}$ this, via an induction argument using the
long exact cohomology sequence, implies the desired vanishing for it as well.
\end{proof}

\end{document}
